\chapter{Policy Implications of School Board Votes}
\label{chap:policy}
\chapterstyle{deposit}



\section{Introduction}

Prior chapters have focused on the ability of community members to exert 
democratic influence on their local school boards. While the evidence on the 
responsiveness of the community to the democratic election of school boards has 
been mixed, it remains to be seen if successfully defeating an incumbent results 
in observable changes in outcomes. Mobilizing to vote or run for school board can 
be expensive relative to the small perceived benefit of doing so. Are voters and 
potential candidates correct if they assess school boards as unable to shift 
policy toward their preferences? Are school boards so constrained as actors that 
they are not able to effect meaningful policy change? In this chapter I explore 
this for two specific policy changes -- improvements in student outcomes and 
the change in district administrator.

The challenge in testing school board responsiveness is great. First, in most 
cases school boards are non-partisan offices and education issues do not fall 
along traditional ideological lines \citep[for counter example:][]{RefWorks:347}. 
This poses a measurement challenge for researchers seeking evidence of policy 
responsiveness. Additionally, the majority of school board election activity is 
low-impact and low-cost, eliminating traditional ways of measuring the information 
provided to voters about candidates and their policy preferences 
\citep{RefWorks:322,RefWorks:321}. Finally, the link between school board decisions 
and aggregate student outcomes in the community is weak in two distinct ways. First, 
the ability or desire of boards to make decisions that are counter to the preferences 
of the district administrator is questionable. % CITATION
Second, school board decisions, like other policy decisions, may take time to 
show results in district outcomes, particularly aggregate measures of student 
performance. 

Despite these measurement difficulties, there are strong reasons to believe that 
school boards do in fact matter. School boards have considerable power over 
local issues in their community. School boards in Wisconsin are responsible for 
setting property tax levies, labor relations with teachers and support staff within 
the school, and often serve as the largest single employer in their community. They 
determine graduation requirements and discipline and attendance policies for the 
students in their schools. The driving idea behind locally elected school boards 
was to ensure that the local community retained some control over the direction 
taken by their schools and that the school system was responsive to their 
preferences. Until now there have only been limited tests of school board 
responsiveness \citep{RefWorks:345}. 

In this chapter I explore the impact of incumbent defeat on student outcomes both 
directly and indirectly in Wisconsin school districts. I focus on incumbent defeat 
because it is the clearest sign of dissatisfaction with current board policy and 
practice and serves as a cornerstone for dissatisfaction theory, the richest 
theory linking school board practice to educational outcomes \citep{RefWorks:297}. 



\section{Literature Review and Theory}

While previous chapters examined the exercise of democratic control 
of school board elections -- under what conditions did voters and candidates 
participate, and was this participation sensitive to policy pressures -- this 
chapter seeks to understand the result of such participation. In order for 
democratic institutions to matter, the results of elections must also matter. 
Challengers must represent a policy alternative to defeated incumbents. Policy 
distinction must be seen among candidates. 

Incumbent defeat provides a strong sign that citizens are dissatisfied with the 
existing outputs of their school government and wish to see changes. Without 
polling data or an in-depth study of the politics of each community, it is 
difficult to understand what might be the source of this voter dissatisfaction 
and what voters might wish to see change. In this study I look for evidence of 
change along the two dimensions previously identified in the literature as 
critical to voters -- the district administrator and student performance. If 
voters are dissatisfied with their board and replace some board members with new 
challengers, it is likely that they demand higher student performance, new 
district management, or both. 

Much of the work referred to in Chapter \ref{chap:intro} dealt with the impact 
of school boards on policy change. In fact, both the dissatisfaction theory and 
the public choice theories of school board politics include superintendent
turnover as key evidence \citep{RefWorks:297,RefWorks:271, RefWorks:351,
RefWorks:339,RefWorks:340,RefWorks:243}. Recently a key improvement in empirical 
tests of this work has focused on the importance of political turnover among 
school board members, as opposed to the more typical turnover due to retirement, 
as a precursor to improved student performance \citep{RefWorks:311}. 

\citet{RefWorks:347} represents the most thorough study of school boards as agents 
of change in social issues. Some view school boards as potentially captured by 
special interest groups \citep{RefWorks:320}. There is evidence that when 
boards are elected, they become more responsive to interest groups than previously 
when they were appointed \citep{RefWorks:242}. 

Boards can influence student outcomes through a range of decisions - some direct 
and some less direct. \citet{RefWorks:210} provides an example by examining the 
relationship between board decisions to consolidate school districts and the 
resulting changes in student achievement. On a more universal basis, 
\citet{RefWorks:214} explores if boards change policies in response to public
accountability for student achivement, and if this results in improvements in
student achievement.\footnote{This argument assumes a certain level of 
effectiveness that is not captured by any measures in the current study and 
may not be evenly distributed across boards \citep{RefWorks:237,RefWorks:261}.} 
In many cases, boards are criticized as being highly constrained actors with 
little room to set policy after complying with state and federal mandates 
\citep{RefWorks:316}, although the extent of this is not clear 
\citep{RefWorks:223}.\footnote{Note that this argument can be reduced down ever
further to the impact of district policies on schools or school policies on
classroom \citep{RefWorks:255}}

Yet boards allocate resources, and resources matter \citep{RefWorks:183}. 

\section{Data}

To study the impacts of incumbent defeat in Wisconsin school districts I employ 
a panel data set of school board election outcomes matched to school district 
administrative records. School board election results were obtained for a sample 
of Wisconsin school districts for the time period from 2002-2012 and coded to 
indicate winners, losers, incumbents, and repeat candidates. To code incumbent 
defeat, I matched a roster of known school board members in a given school year 
against the victors in the following school year to identify which candidates 
were incumbents. 

This method differs from the board and candidate self-reported measure of incumbent 
defeat that have been used previously in the literature \citep{RefWorks:345}. The 
longitudinal design is also more observations per district and includes a larger 
proportion of districts in the state than prior designs. In addition to this, 
a host of administrative data on school districts is also included, most critically 
measures of student outcomes and superintendent turnover. I discuss these in 
turn below. 

\subsection{Independent Variables}

Previous chapters have established that school board incumbent defeat is rare 
because incumbents are rarely challenged, and even more rarely is that challenge 
serious. Previous research on school board incumbent defeat has focused on 
correctly measuring political and apolitical turnover by focusing on incumbents 
who are defeated, not simply those who do not return to the board for another term. 

Timing is of critical importance in assessing the \emph{impact} of school board 
incumbent defeat. While school board members are elected in spring elections, 
their term begins the following fall. I use a lagged indicator of incumbent defeat 
to account for this by looking at the impact incumbent defeat has on student 
performance on the state exam and graduation rate two years after the election date. 
For an incumbent defeated in the spring election of 2004, the impact of this 
defeat is expected to come not during the 2005-06 school year, but instead the 
2006-07 year.\footnote{I check models with both a twice and single lagged measure 
of incumbent defeat for substantive differences.}

\begin{knitrout}
\definecolor{shadecolor}{rgb}{0.969, 0.969, 0.969}\color{fgcolor}\begin{figure}
\includegraphics[width=\maxwidth]{../../phd_figs/chap6-incDefeatMeasure-1} \caption[Incumbent Defeat Year to Year in Wisconsin 2002-2012]{Incumbent Defeat Year to Year in Wisconsin 2002-2012}\label{fig:incDefeatMeasure}
\end{figure}


\end{knitrout}

Figure \ref{fig:incDefeatMeasure} shows the count and proportion of  
incumbent defeats in the Wisconsin data from 2002-2012. Typically 40-60 incumbents 
are defeated a year, representing approximately 10\% of all incumbents running 
each year. One shortcoming with using simple incumbent defeat as a measure is that 
Wisconsin school boards are majority rule decision making bodies meaning that if 
a single incumbent is defeated, it may be insufficient to result in substantive 
policy change depending on the policy positions of other members of the board. 
Figure \ref{fig:incTurnover} shows the number of seats on average up for election 
in districts in each year. The majority of districts elect one or two members 
and this represents a minority of the typically sized board of five to seven 
members. 

\begin{knitrout}
\definecolor{shadecolor}{rgb}{0.969, 0.969, 0.969}\color{fgcolor}\begin{figure}
\includegraphics[width=\maxwidth]{../../phd_figs/chap6-incTurnover-1} \caption[Number of Board Seats Elected Each Year]{Number of Board Seats Elected Each Year}\label{fig:incTurnover}
\end{figure}


\end{knitrout}

Unfortunately, there is not enough variation in the number of board seats elected 
and the number of incumbents defeated to use anything other than a binary measure 
of incumbent defeat. 

\subsection{Control Variables}

Following the work of \citet{RefWorks:345}, I include relatively few control 
variables. Including too many control variables results in overly restricting the 
variance available to be explained by incumbent defeat and possibly collinear. 
However I do include measures of whether the district is rural, urban, or exurban, 
and a control for the size of the voting age population. I also include measures 
of the level of education of the adult population and the percentage of adults in 
the community who are over the age of 65. I also include a measure of economic 
disadvantage for students in the school district -- a key predictor of student 
outcomes. 

\subsection{Outcomes}

First, and most directly, boards are charged with hiring and monitoring 
the performance of a district administrator who manages the academic and 
financial affairs of the district. The choice of superintendent is critical to 
the board as the superintendent is the board's primary conduit of information 
about the health of the district and the educational strategies in place for 
students \citep{RefWorks:271}. Superintendent turnover is a primary prediction 
of dissatisfaction theory, and rightfully so. If voters are unhappy with the 
state of affairs in their school district, in all likelihood they are unhappy 
with the performance of their district administrator \citep{RefWorks:246,RefWorks:245,
RefWorks:273}. 

Superintendent change is coded using the Wisconsin Department of Public Instruction 
database on school staffing - the PI-1202 report. This report includes annual 
information on the job position, license status, full-time equivalency or full-time 
part-time nature of the job, and salary. This information is collected each fall. 
A superintendent change was coded if a new superintendent was listed that differed 
from the prior year. 

Figure \ref{fig:supChange} shows the percentage of school districts in the sample 
each year that experienced a superintendent turnover. Annually around 15\% of 
districts in the sample had a change in superintendent. 

\begin{knitrout}
\definecolor{shadecolor}{rgb}{0.969, 0.969, 0.969}\color{fgcolor}\begin{figure}
\includegraphics[width=\maxwidth]{../../phd_figs/chap6-supChange-1} \caption[Superintendent turnover by year in sample]{Superintendent turnover by year in sample.}\label{fig:supChange}
\end{figure}


\end{knitrout}


More indirectly, school boards are responsible for setting policy and directing 
the superintendent to provide quality education. Thus, incumbent defeat may 
serve as a method by which the community can express its demand for higher 
quality education as measured by accountability measures such as student 
test scores and graduation rates. While prior research has found evidence of 
an effect of school board turnover and elections on student performance on 
standardized tests, this outcome seems the furthest from the control of school 
board members. 



School boards appear to have more ability to influence graduation rates. First, 
school boards in Wisconsin have wide leverage in setting the academic standards 
for graduation and there is large variability in the credit requirements to graduate 
in high schools across the state. Second, graduation rates are highly visible 
accountability measures throughout this period of observation, and districts may 
have been able to attempt a wide variety of policy changes to improve graduation 
outcomes relatively quickly. 

The graduation rate measure I use is the only measure consistently calculated 
throughout the time period and the one used for accountability purposes for 
most of the period from 2002-2012 -- the legacy graduation rate. This rate differs 
from more recent graduation rates reported by schools, districts, and states -- 
it is not a four-year graduation rate. One challenge with this measure is that 
overall Wisconsin has a high graduation rate, and many school districts in 
Wisconsin has particularly high graduation rates. Figure \ref{fig:gradRateplot} 
shows the graduation rates of school districts throughout the period of study. 

\begin{knitrout}
\definecolor{shadecolor}{rgb}{0.969, 0.969, 0.969}\color{fgcolor}\begin{figure}
\includegraphics[width=\maxwidth]{../../phd_figs/chap6-gradRateplot-1} \caption[Public High School Graduation Rates]{Public High School Graduation Rates}\label{fig:gradRateplot}
\end{figure}


\end{knitrout}


\section{Methods}

Measuring the impact of school board turnout on superintendent turnover and 
student outcomes is difficult because of the large number of competing 
forces shaping these outcomes. Student performance is largely predicted by 
family income and parental education levels. Superintendent turnover may be 
related to salary and the prestige associated with leading certain school 
districts. In order to better understand the impact of school board turnover 
net of these factors, I employ a regression strategy to control for as many 
of these observable and unobservable factors as possible. 


I employ two distinct modeling strategies. The first is the most restrictive. 
As in previous chapters, I employ the bounding strategy suggested by \citet{mhecon} 
to leverage the repeated observations in my data to compare, where possible, 
school districts within themselves. The disadvantage of this strategy is that 
for the outcomes in question, relatively little variance remains to be explained 
after accounting for all between school district effects, because the within 
school district variation is so small. 

The second strategy is more susceptible to unobserved alternative explanatory 
variables, but attempts to explain the variation in outcomes between school 
districts by comparing school districts with similar demographic 
charcteristics, but which differ on their incumbent defeat. 

\section{Results}


\subsection{Student Outcomes}

\begin{table}[!ht]
\caption{Incumbent Defeat and Reading Proficiency Rates}
\label{tab:readMod} 
\begin{tabular}{ l D{.}{.}{3}D{.}{.}{3}D{.}{.}{3} } 
\hline 
  & \multicolumn{ 1 }{ c }{ Simple } & \multicolumn{ 1 }{ c }{ Lag } & \multicolumn{ 1 }{ c }{ FE } \\ \hline
 %                 & Simple            & Lag               & FE               \\ 
Intercept         & 3.609 ^{***}      & 0.758 ^{***}      & 0.770            \\ 
                  & (0.096)           & (0.101)           & (1.367)          \\ 
Inc. Def. Lag     & -0.024 ^\dagger  & 0.007             & 0.001            \\ 
                  & (0.013)           & (0.007)           & (0.007)          \\ 
Exurb             & 0.104             & -0.003            & 1.032 ^{***}     \\ 
                  & (0.064)           & (0.012)           & (0.229)          \\ 
Rural             & 0.069             & 0.005             & 1.029 ^{***}     \\ 
                  & (0.053)           & (0.012)           & (0.229)          \\ 
Board Size        & -0.014            & 0.000             & -0.009           \\ 
                  & (0.011)           & (0.003)           & (0.009)          \\ 
FRL \%           &                   & -0.002 ^{***}     & 0.001            \\ 
                  &                   & (0.000)           & (0.001)          \\ 
Prof Rate Lag     &                   & 0.797 ^{***}      &                  \\ 
                  &                   & (0.027)           &                  \\ 
VAP               &                   & -0.002            & 0.245 ^\dagger  \\ 
                  &                   & (0.004)           & (0.130)          \\ 
\% Bach. +       &                   & 0.201 ^*          & 0.079            \\ 
                  &                   & (0.085)           & (0.535)          \\ 
\% Over 65       &                   & 0.103             & -0.069           \\ 
                  &                   & (0.096)           & (1.106)           \\
 $N$               & 1918              & 1425              & 1918             \\ 
$R^2$             & 0.017             & 0.846             & 0.906            \\ 
adj. $R^2$        & 0.015             & 0.845             & 0.888            \\ 
Resid. sd         & 0.235             & 0.094             & 0.079             \\ \hline
 \multicolumn{4}{l}{\footnotesize{$^\dagger$ significant at $p<.10$; $^* p<.05$; $^{**} p<.01$; $^{***} p<.001$}}\\
\multicolumn{4}{l}{\footnotesize{District fixed effects not displayed.}} 
\end{tabular} 
 \end{table}




 
\begin{tabular}{ l D{.}{.}{3}D{.}{.}{3}D{.}{.}{3} } 
\hline 
  & \multicolumn{ 1 }{ c }{ Lag } & \multicolumn{ 1 }{ c }{ Model 2 } & \multicolumn{ 1 }{ c }{ Model 3 } \\ \hline
 %                & Lag              & Model 2          & Model 3         \\ 
Intercept        & 3.787 ^{***}     & 0.907 ^{***}     & -5.004 ^*       \\ 
                 & (0.099)          & (0.109)          & (2.132)         \\ 
Inc. Def. Lag    & -0.023           & 0.007            & -0.004          \\ 
                 & (0.015)          & (0.007)          & (0.007)         \\ 
Exurb            & 0.112 ^\dagger  & 0.037 ^*         & 1.395 ^{***}    \\ 
                 & (0.064)          & (0.017)          & (0.320)         \\ 
Rural            & 0.068            & 0.033 ^*         & 1.391 ^{***}    \\ 
                 & (0.053)          & (0.016)          & (0.320)         \\ 
Board Size       & -0.006           & -0.000           & -0.007          \\ 
                 & (0.011)          & (0.003)          & (0.011)         \\ 
FRL \%          &                  & -0.002 ^{***}    & 0.003 ^*        \\ 
                 &                  & (0.000)          & (0.001)         \\ 
Prof Rate Lag    &                  & 0.746 ^{***}     &                 \\ 
                 &                  & (0.025)          &                 \\ 
VAP              &                  & 0.011 ^*         & 0.775 ^{***}    \\ 
                 &                  & (0.005)          & (0.204)         \\ 
\% Bach. +      &                  & 0.159            & 1.117 ^\dagger \\ 
                 &                  & (0.102)          & (0.592)         \\ 
\% Over 65      &                  & 0.052            & 0.543           \\ 
                 &                  & (0.086)          & (1.265)          \\
 $N$              & 1918             & 1425             & 1918            \\ 
$R^2$            & 0.010            & 0.840            & 0.889           \\ 
adj. $R^2$       & 0.008            & 0.839            & 0.868           \\ 
Resid. sd        & 0.250            & 0.097            & 0.091            \\ \hline
 \multicolumn{4}{l}{\footnotesize{$^\dagger$ significant at $p<.10$; $^* p<.05$; $^{**} p<.01$; $^{***} p<.001$}}\\
\multicolumn{4}{l}{\footnotesize{District fixed effects not displayed.}} 
\end{tabular} 


% HSC 

\begin{table}[!ht]
\caption{Incumbent Defeat and High School Graduation}
\label{tab:gradMod} 
\begin{tabular}{ l D{.}{.}{3}D{.}{.}{3}D{.}{.}{3} } 
\hline 
  & \multicolumn{ 1 }{ c }{ Simple } & \multicolumn{ 1 }{ c }{ Lag } & \multicolumn{ 1 }{ c }{ FE } \\ \hline
 %                 & Simple            & Lag               & FE               \\ 
Intercept         & 4.501 ^{***}      & 2.525 ^{***}      & 3.790 ^{***}     \\ 
                  & (0.024)           & (0.269)           & (0.629)          \\ 
Inc. Def. Lag     & -0.005 ^\dagger  & -0.002            & -0.002           \\ 
                  & (0.003)           & (0.003)           & (0.003)          \\ 
Exurb             & 0.051 ^*          & 0.016             & 0.092            \\ 
                  & (0.022)           & (0.011)           & (0.112)          \\ 
Rural             & 0.059 ^{**}       & 0.021 ^\dagger   & 0.105            \\ 
                  & (0.021)           & (0.010)           & (0.109)          \\ 
Board Size        & -0.002            & 0.001             & 0.009            \\ 
                  & (0.002)           & (0.001)           & (0.007)          \\ 
FRL \%           &                   & 0.452 ^{***}      &                  \\ 
                  &                   & (0.058)           &                  \\ 
Grad Rate Lag     &                   & -0.008 ^{**}      & 0.053            \\ 
                  &                   & (0.003)           & (0.059)          \\ 
VAP               &                   & 0.150 ^{***}      & 0.212            \\ 
                  &                   & (0.032)           & (0.208)          \\ 
\% Bach. +       &                   & -0.073 ^\dagger  & 0.393            \\ 
                  &                   & (0.037)           & (0.400)          \\ 
\% Over 65       &                   &                   & 0.001 ^\dagger  \\ 
                  &                   &                   & (0.000)           \\
 $N$               & 1755              & 1755              & 1755             \\ 
$R^2$             & 0.059             & 0.311             & 0.592            \\ 
adj. $R^2$        & 0.056             & 0.307             & 0.512            \\ 
Resid. sd         & 0.059             & 0.051             & 0.043             \\ \hline
 \multicolumn{4}{l}{\footnotesize{$^\dagger$ significant at $p<.10$; $^* p<.05$; $^{**} p<.01$; $^{***} p<.001$}}\\
\multicolumn{4}{l}{\footnotesize{District fixed effects not displayed.}} 
\end{tabular} 
 \end{table}




\subsection{Superintendent Turnover}

 
\begin{tabular}{ l D{.}{.}{3}D{.}{.}{3}D{.}{.}{3} } 
\hline 
  & \multicolumn{ 1 }{ c }{ Simple } & \multicolumn{ 1 }{ c }{ Full } & \multicolumn{ 1 }{ c }{ Fixed Effect } \\ \hline
 %                         & Simple           & Full             & Fixed Effect    \\ 
(Intercept)               & -1.455 ^{***}    & -1.245           & -1.807 ^*       \\ 
                          & (0.268)          & (0.936)          & (0.721)         \\ 
incsDefeatedFac\_lag1    & 0.271 ^\dagger  & 0.244            & 0.107           \\ 
                          & (0.159)          & (0.159)          & (0.204)         \\ 
locale2\_lagexurb        & -0.462           & -0.317           & -0.299          \\ 
                          & (0.319)          & (0.356)          & (0.280)         \\ 
locale2\_lagrural        & -0.329           & -0.176           & -0.075          \\ 
                          & (0.273)          & (0.351)          & (0.221)         \\ 
econ\_disadv\_per\_lag &                  & 0.008 ^\dagger  & 0.021 ^*        \\ 
                          &                  & (0.005)          & (0.008)         \\ 
log(VAP\_adj\_lag)      &                  & 0.044            & 0.030           \\ 
                          &                  & (0.080)          & (0.050)         \\ 
PerBachelorOrAbove\_lag  &                  & -2.801 ^*        & -1.458          \\ 
                          &                  & (1.308)          & (1.450)         \\ 
per65o\_lag              &                  & -3.797 ^*        & -5.599 ^{**}    \\ 
                          &                  & (1.840)          & (1.890)          \\
 $N$                       & 1918             & 1918             & 1918            \\ 
AIC                       & 1632.418         & 1626.880         & 1946.831        \\ 
BIC                       & 1721.363         & 1804.769         & 8973.456        \\ 
$\log L$                 & -800.209         & -781.440         & 290.584          \\ \hline
 \multicolumn{4}{l}{\footnotesize{$^\dagger$ significant at $p<.10$; $^* p<.05$; $^{**} p<.01$; $^{***} p<.001$}} 
\end{tabular} 



\clearpage

\section{Conclusion}

This chapter has looked for evidence of the impact of school board electoral 
defeat on student outcomes and on school district leadership. 



